\section{Method overview}
\label{sec:overview}

%%----------------------------------------------------------------
%% Overview Figure
%%----------------------------------------------------------------
\begin{figure*}[t]
	\begin{center}
	\includegraphics[width=0.9\linewidth]{./images/overview-hd.png}
% 	\includegraphics[height=4cm]{./images/patches-empty-patch-with-pattern.png}
% 	\includegraphics[height=4cm]{./images/patches-connected-patch.png}
% 	\includegraphics[height=4cm]{./images/patches-collision-free-patch.png}
% 	\includegraphics[height=4cm]{./images/patches-smoothed-patch.png}
	\caption{
		\textbf{Overview} Input and output points in a patch's patterns are initially connected and subsequently modified and smoothed out to be collision free.
	}
	\label{fig:overview}
	\end{center}
\end{figure*}
%%----------------------------------------------------------------
%% Overview Figure
%%----------------------------------------------------------------

\note{Overview: remind definitions on crowd patches: patch, pattern, spatiotemporal waypoints (boundary conditions: input, output, initial states // boundary conditions should be strictly enforced. Movable control points), period of time \dots}

\note{Starting Definitions:}

\begin{description}

\item[Patch - ]{A patch is a set $\{ A, \pi, D, S\}$ where $A$ is a geometrical area with a convex polygonal shape, $\pi$ the period of time of the animation, $D$ and $S$ are the sets of dynamic and static objects, respectively. These last two sets may be empty in case of an empty patch.}

\item[Static obstacles -]{Static objects are simple obstacles whose geometry is fully contained inside the patch.}

\item[Dynamic objects]{Dynamic objects are animated: they are moving in time according to a set of trajectories $T$.}

\end{description}

We define a {\bf trajectory} inside a patch as a function going from time to position, more specifically from the subset $[ 0,\pi ]$ to $A$:$ \tau:[t_1,t_2]\rightarrow A, \hspace{0.3cm} 0 \leq t_1 < t_2 \leq \pi$. We represent a trajectory as a list of control points connected by segments. 

A {\bf control point} is a point in space and time. All control points in a trajectory will either be a movable control point or a boundary control point. Boundary control points serve as entry and exit points to the patch and cannot be moved, added or deleted. Movable control points can be moved, added, or removed from the trajectory as long as they do not violate the \note{constraints of the patch, their position must be inside $A$ and the time between $t_1$ and $t_2$.} \\
\devin{-- Should we state exactly whate these constraints are here? }\\
\memo{-- Sure, I added a line for that}

A {\bf segment} is a straight line connecting two control points in a specific order. Since these are unidirectional lines in space-time, it is important to remember that they are not allowed to go backwards in time.

There are two categories of dynamic objects: endogenous agents and exogenous agents. {\bf Endogenous agents} remain inside $A$ for the total period of time $\pi$. In order to achieve periodicity for the animation, they are associated with a trajectory $\tau : [0,\pi] \leftarrow A$, such that it respects the periodicity condition: the position at the start and at the end of the animation must be the same, i.e. \mbox{$\tau (0) = \tau (\pi)$}.\\
{\bf Exogenous agents} go outside $A$. They enter the patch at time $t_{initial}$ and position $a_{initial}$, and they exit at time $t_{final}$ and position $a_{final}$. For each agent we associate a sequence of $n$ trajectories $\{ \tau_1, \tau_2, \dots, \tau_n\}$. Sequences may have only one trajectory, but some agents require additional trajectories in order to satisfy speed and time constraints. The following conditions must be respected in each sequence of trajectories associated with an exogenous agent:

\begin{itemize}

\item{$a_{initial}$ and $a_{final}$ must be points in the border of $A$. Otherwise, they couldn't be exogenous agents.}

\item{If the sequence is composed by more than one trajectory, for each two contiguous trajectories, the following must be true to ensure continuity: $\tau_i(\pi) = \tau_{i_{next}}(0)$.}

\end{itemize}

Note that the second condition implicitly implies that in sequences with multiple trajectories, each middle trajectory must be fully defined in the period of time $[0,\pi]$, while $\tau_1$ must be defined in $[t_{initial},\pi]$ and $\tau_n$ must be defined in $[0, t_{final}]$.

\begin{description}

\item[Pattern -]{If a patch is a spatio-temporal cube (or any other right prism, depending on the type of polygon used as its area, then a pattern could be defined as one lateral side of the \note{cube (or right prism).} Specifically, it is a rectangle whose base is one of the edges of the polygonal area (we define $I$ as this two dimensional vector), and its height is equal to the period. These patterns also include the sets of boundary control points. We divide this set evenly into two subsets: Input and Output. The Input set contains the boundary control points where exogenous agents begin their trajectories; we call these Entry Points. Conversely the elements of Output are called Exit Points. They establish the position in time and space that the exogenous agents finish their paths. Formally defined, a patch P is:
$$ P = {l, p_i, I_i [p_i, t_i], O_j[p_j, t_j]}$$
We populate virtual environments by sticking patches together. Thus, we have to ensure continuity between trajectories for exogenous agents passing through two contiguous patches. This means that two adjacent patches must \note{have a similar pattern on the side they share}. The vector $l$ \claudia{Is this vector $l$ or $I$?, please check } and period must be equivalent and the sets of Input and Output are exchanged. If we have $P_1={l, pi, I [p_1,t_1], O[p_2,t_2]}$ and $P_2={l�, pi�, I�[p�_2, t�_2], O�[p�_1, t�_1]}$, then, in order to satisfy $C^0$ continuity we must ensure: 
$$ l=l�, pi=pi�, p_1=p�_1, t_1=t�_1, p_2=p�_2, t_2=t�_2.$$
We then say $P_1$ is the mirror pattern of $P_2$. In the animation, this will be seen as an agent going from one patch to an adjacent one. 
If the area of a patch is a square, the patch defines 4 patterns, one for each of its sides. Patterns defined by a patch have the property that the sum of the cardinality of all the Inputs is the same as the sum of the cardinality of all Outputs. We call this the parity condition: $\sum(|Inputs|)= \sum(|Outputs|)$.

A patch defines a set of Patterns, and conversely, a set of patterns satisfying the parity condition, having the same period, and whose vectors define a convez polygonal area, can be used to create a patch. 
 }

\end{description}



